%% ============================================================
%% 片段：s01-intro   对应原文 PDF 第 2--5 页
%% 本文件由 gw/convert.py 从 arXiv 源码（hep-th/0612073v2）自动生成：
%%   * 公式/编号/引用/图 已经锁定为原文内容，请勿改动任何数学内容与编号；
%%   * 只需把英文正文译为中文，并在每节开头加“本节重点”、关键处加译注。
%% ============================================================


\section{引言}

%% ==== 原文 §1：PDF 第 2 页（印刷第 1 页）====
\label{sec:1}

\begin{keytakeaways}
  \begin{itemize}[leftmargin=1.4em,itemsep=2pt,topsep=2pt]
    \item \textbf{物理图像}：四维 ${\mathcal N}=4$ 超杨–米尔斯理论经扭曲、在黎曼面上紧化后，低能约化为靶空间是希格斯丛的希钦模空间的西格玛模型；朗兰兹纲领中的分歧，在规范场论中对应沿余维二子流形引入的面算符，而非支撑在一维流形上的威尔逊算符或 't Hooft 算符。
    \item \textbf{数学工具}：具有单极点的带分歧希格斯丛（即带抛物结构的希格斯丛）；仿射外尔群（affine Weyl group）在其模空间上同调上的作用，以及（第~\ref{sec:4}~节）仿射辫群在 $A$-膜、$B$-膜上的作用；这些作用与施普林格表示、仿射赫克代数的 Kazhdan--Lusztig 理论、Bezrukavnikov 关于导出范畴的作用等表示论结构相类比。
    \item \textbf{关键结论}：对 $S$-对偶在面算符参数上的作用给出自然提议（第~\ref{sec:2}~节），并据此提出驯服分歧情形下几何朗兰兹纲领应有的表述（第~\ref{sec:4}~节）；分歧点附近的局部行为由靶空间为复余伴随轨道的西格玛模型描述（第~\ref{sec:6}~节）。
    \item \textbf{边界与展望}：本文只处理驯服分歧，野分歧留待后续工作 \cite{ref33}；共形场论给出的、允许任意阶极点的统一处理 \cite{ref6,ref7} 则未被涉及。
  \end{itemize}
\end{keytakeaways}

数论中的朗兰兹纲领（Langlands program）\cite{ref1} 把数域（number field）的伽罗瓦群（Galois group）的表示与自守形式（automorphic form，例如 $SL(2,\mathbb{Z})$ 的普通模形式）联系起来。它也有一个几何类比，即以普通黎曼面（Riemann surface）代替数域。结果表明，这一几何类比与量子场论（quantum field theory）密切相关。人们曾借助二维共形场论（conformal field theory）\cite{ref2,ref3,ref4,ref5,ref6} 对它作了广泛研究，近来又通过四维规范场论（gauge theory）与电磁对偶（electric-magnetic duality）\cite{ref8} 来研究它。更多的解释与文献可参见 \cite{ref8} 的引言以及近期的一篇综述 \cite{ref9}。

几何朗兰兹对应（geometric Langlands correspondence）的最简单形式，一侧是黎曼面 $C$ 上的一个平坦联络（flat connection），另一侧是称为 ${\mathcal D}$-模（${\mathcal D}$-module）的一种更精致的结构。该问题有一个推广：从 $C$ 中删去有限多个点 $p_1,p_2,\dots,p_n$，并以 $C'=C\backslash\{p_1,\dots,p_n\}$ 上的一个平坦联络为出发点，该联络在给定点附近具有指定的奇点（singularity）。这种情形给出数论中所谓分歧（ramification）的几何类比。由于分歧在数论中几乎无法回避，把几何朗兰兹纲领推广到分歧情形，是数论与几何之间这一类比的重要组成部分。该问题还有一个重要的局部版本，着眼于分歧点附近的行为。

\begin{yizhu}[含糊点拨]
“局部版本”指只看单个分歧点 $p$ 附近的行为：把 $C$ 换成以 $p$ 为圆心的圆盘（或去心圆盘），只要求该圆盘上的联络在 $p$ 附近有指定奇点，而不必给出整条 $C$ 上的联络。两个版本在本文中分工不同：第~\ref{sec:4}~节的提议是全局的，靶空间是整条 $C$ 的希钦模空间 $\MH$；直到第~\ref{sec:6}~节（PDF 第 130 页起）才给出局部描述，靶空间换成复余伴随轨道。第~\ref{sec:3.7}~节末尾（PDF 第 68 页起）本文明言：若能在 $\MH$ 的框架里看到仿射赫克代数，将改进“全局问题（$\MH$）”与“局部问题（复余伴随轨道）”之间的类比——所谓“局部”，正指后者那种只依赖分歧点附近数据的问题。
\end{yizhu}

本文的目标是把几何朗兰兹纲领的规范场论方法推广到涵盖分歧情形。基本想法是：在朗兰兹纲领中允许分歧，在规范场论中对应于引入面算符（surface operator）；这类算符在某种程度上类似于威尔逊算符或 't Hooft 算符，但它支撑在二维流形（two-manifold）上，而非一维流形（one-manifold）上。相关的面算符通过在余维二子流形（codimension two submanifold）上指定某类奇点来定义。规范场论中大致属于这一类型的余维二奇点，已在多种背景下被研究过，包括宇宙弦（cosmic string）理论 \cite{ref10}、Donaldson 理论 \cite{ref11}、拓扑场论（topological field theory，见 \cite{ref12} 第 5 节）、规范场论的动力学与黑洞 \cite{ref13,ref14,ref15}、瞬子（instanton）与塞伯格--威滕理论及可积系统的关系 \cite{ref16,ref17}，以及弦理论：在弦理论中，本文所考虑算符的一些特例可通过相交膜（intersecting brane）构造出来 \cite{ref18}，我们将在第~\ref{sec:6.4}~节中加以描述。

\begin{yizhu}[易错点]
面算符并非只由奇异性的类型决定，它带有连续参数。具体地，在面算符所沿的二维子流形 $D$ 的法平面上取极坐标 $x^1+ix^2=re^{i\theta}$，并要求 $A=\alpha\,d\theta$、$\phi=\beta\,dr/r-\gamma\,d\theta$（第~\ref{sec:2.1}~节，PDF 第 5--6 页）。其中 $\alpha$ 只在模去余特征格 $\Lambda_{\rm cochar}$ 的意义下有定义：绕 $D$ 一周的和乐为 $\exp(-2\pi\alpha)$，此处的 $2\pi$ 因子与负号来自本文的约定 $P\exp(-\int A)$；三元组 $(\alpha,\beta,\gamma)$ 还要再模去外尔群的作用。$\beta$ 与 $\gamma$ 则是希格斯场特征多项式中极点的系数，与 $\alpha$ 扮演不同的角色。
\end{yizhu}

\begin{yizhu}[物理图像补充]
面算符与威尔逊算符、\,'t Hooft 算符的差别首先在支撑维数：四维理论中支撑在 1 维曲线上的算符余维为 3，支撑在 2 维曲面上的算符余维为 2。线算符的标签离散——威尔逊算符由 $G$ 的表示标定，\,'t Hooft 算符由余权（等价地，对偶群的表示）标定；面算符则带连续参数。据第~\ref{sec:2.1}~节（PDF 第 5 页起），沿曲面 $D$ 的奇点由三元组 $(\alpha,\beta,\gamma)\in\TT\times\mathfrak t\times\mathfrak t$ 刻画并模去外尔群；$D$ 的法平面是复一维的，其上 $A=\alpha\,d\theta$ 的系数可取环面 $\TT$ 中的连续值，这正是面算符能承载朗兰兹侧分歧数据的原因。
\end{yizhu}

下面我们简要说明这些规范场论奇点与希格斯丛（Higgs bundle）理论之间的联系。几何朗兰兹的规范场论方法基于经扭曲（twist）、并在黎曼面上紧化的 ${\mathcal N}=4$ 超杨–米尔斯理论（super Yang-Mills theory）。如此紧化后的理论在低能下 \cite{ref19,ref20} 约化为一个西格玛模型（sigma model），其靶空间是一个超 Kähler（hyper-Kähler）流形，即希格斯丛的希钦模空间（Hitchin moduli space）\cite{ref21}。这一图像能够容纳何种余维二奇点？在最基本的情形，即 $C$ 上的平坦联络只有单极点（simple pole）时，相应的奇点已由 Simpson \cite{ref22} 给出描述与分析。具有这类奇点的希格斯丛，就是我们所谓的带分歧的希格斯丛（ramified Higgs bundle）。（它们也被称为带抛物结构（parabolic structure）的希格斯丛。）与之相伴的超 Kähler 模空间已由 Konno \cite{ref23} 构造，其拓扑已由 Nakajima \cite{ref24} 阐明。对于具有高阶极点的希格斯丛，相应的理论已由 Biquard 与 Boalch \cite{ref25} 发展起来。单极点与高阶极点这两种情形，分别对应于朗兰兹纲领语境中所谓的驯服分歧（tame ramification）\footnote{有时，“驯服分歧”一词被更狭窄地用来指具有幂幺单值性（unipotent monodromy）的平坦丛的情形。我们不作这一限制，而考虑具有任意单极点的联络与希格斯丛，如式~\eqref{eq:2.2} 与第~\ref{sec:3.3}~节所述。}与野分歧（wild ramification）。本文集中讨论驯服分歧。

\begin{yizhu}[物理图像补充]
“扭曲并紧化”省略了两步：先对 ${\mathcal N}=4$ 超杨–米尔斯理论作 GL 扭曲（其参数为 $t$，见 \cite{ref8}），再把四维流形取为 $C$ 与另一个二维流形的乘积，并沿 $C$ 紧化。$t$ 在约化到二维后消失，剩下的恰是希钦方程，因此低能有效理论是以希钦模空间为靶空间的西格玛模型。就几何朗兰兹而言，四维的 $S$-对偶在二维侧表现为镜像对称，面算符则表现为膜；这正是第~\ref{sec:4}~节把几何朗兰兹问题翻译成膜与 ${\mathcal D}$-模语言的出发点。
\end{yizhu}

\begin{yizhu}[含糊点拨]
“带抛物结构的希格斯丛”是同一对象的另一名称：在分歧点 $p_i$ 处，抛物结构指把结构群约化为抛物子群（等价地，给出一面旗）并附上一组权重。本文的单极点情形下，这组数据由面算符参数决定，故两种叫法等价；详细讨论见第~\ref{sec:3.4}~节（PDF 第 48 页起）以及 Simpson \cite{ref22}。此外，脚注中“驯服”的狭义用法只允许幂幺单值性；本文取解析意义，即联络或希格斯场至多只有单极点（正则奇点），残数可任意，这正是面算符参数能够连续变化的根源。
\end{yizhu}

在第~\ref{sec:2}~节中，我们在四维规范场论的语境下考虑对应于单极点的奇点。对 $S$-对偶（$S$-duality）如何作用于参数，我们给出一个自然的提议。第~\ref{sec:3}~节进一步探讨经典几何。该节的亮点是仿射辫群（affine braid group）在带分歧的希格斯丛的模空间的上同调（cohomology）上的作用；第~\ref{sec:4}~节把这一作用推广为仿射辫群在 $A$-膜与 $B$-膜上的作用。这些现象与表示论中发现的若干结构有着密切的亲缘关系，包括外尔群（Weyl group）的施普林格表示（Springer representation）\cite{ref26}、仿射赫克代数（affine Hecke algebra）表示的 Kazhdan--Lusztig 理论 \cite{ref27,ref28}，以及 Bezrukavnikov 新近的工作 \cite{ref29}——把上述理论推广为仿射辫群在施普林格消解（Springer resolution）的导出范畴（derived category）上的作用。关于其中部分材料的阐述，可参见 Chriss 与 Ginzburg 的著作 \cite{ref30}。用超 Kähler 消解（hyper-Kähler resolution）的参数来解释这类结果（正如我们将在希格斯丛模空间的情形中所做的那样），最早由 Atiyah 与 Bielawski \cite{ref31} 在余伴随轨道（coadjoint orbit）与 Slodowy 切片的语境下提出。

\begin{yizhu}[原文疑误]
上一段首句原文作 “In sections 2, we consider...”，用了复数 sections，而实际只指第~\ref{sec:2}~节（PDF 第 4 页，印刷第 3 页），此处应为单数。译文按“第~\ref{sec:2}~节”处理；为不打断原段落，此注置于段末。
\end{yizhu}

\begin{yizhu}[原文疑误]
此处原文称仿射辫群作用在模空间的上同调上，第 4 节再把它推广为仿射辫群在 $A$-膜、$B$-膜上的作用。但按第~\ref{sec:3.7}~节（PDF 第 68 页起）的正文，在参数空间 $\mathfrak t^3$ 模去仿射外尔群 $\AffWeyl$ 之后，上同调作为该商空间上某个平坦丛的纤维而变化，其单值性给出的作用是仿射外尔群在上同调（以及 $K$-理论）上的作用；对膜的作用要到第~\ref{sec:4.5}~节（PDF 第 98 页起）才出现，即该作用被“加厚”为 $A$-膜、$B$-膜范畴上的作用。因此引言此句第一处“仿射辫群”疑为“仿射外尔群”之误。译文照原文译出，建议以正文为准。
\end{yizhu}

\begin{yizhu}[含糊点拨]
所谓“自然的提议”直到第~\ref{sec:2.4}~节才给出具体内容：基本电磁对偶 $S$ 把 $\tau=\theta/2\pi+4\pi i/g^2$ 变为 $-1/(n_{\mathfrak g}\tau)$（$n_{\mathfrak g}$ 对单边（simply-laced）的 $G$ 取 1，否则取 2 或 3），并把 $G$ 换成对偶群（dual group）${}^LG$。面算符参数中，$\beta$ 与 $\gamma$ 作为希格斯场特征多项式的极点系数作相应的简单变换，而 $\alpha$ 与对偶环面上的参数 $\eta$ 按 $(\alpha,\eta)\to(\eta,-\alpha)$ 交换。只在引言层面阅读时，容易误以为对偶仅作用在 $\tau$ 上。
\end{yizhu}

\begin{yizhu}[含糊点拨]
施普林格表示是外尔群在施普林格纤维（施普林格消解在幂零元素处的纤维）的上同调上的自然作用。Atiyah 与 Bielawski \cite{ref31} 建议这样理解它：让余伴随轨道（及其 Slodowy 切片）上的超 Kähler 度量依赖参数，再取参数空间上的单值性，便得到外尔群的作用。第~\ref{sec:3.7}~节（PDF 第 68 页起）把本文结果放进同一机制：$\MH(\alpha,\beta,\gamma;p)$ 的上同调作为 $\mathfrak t^3$ 模去 $\AffWeyl$ 的商上某个平坦丛的纤维而变化，其单值性给出仿射外尔群的作用。引言说“用超 Kähler 消解的参数来解释这类结果”，指的正是这一“参数—单值性—群作用”的机制。
\end{yizhu}

\begin{yizhu}[物理图像补充]
这三条表示论线索构成一个三级阶梯：群作用从外尔群提升到仿射赫克代数、再到仿射辫群，载体从上同调提升到等变 $K$-理论、再到导出范畴。本文的贡献是把整个阶梯实现到规范场论里：第~\ref{sec:3.7}~节先得到仿射外尔群在 $\MH$ 的上同调（以及 $K$-理论）上的作用；第~\ref{sec:4.5}~节（PDF 第 98 页起）把它加厚为仿射辫群在 $A$-膜、$B$-膜范畴上的作用；第~\ref{sec:5}~节（PDF 第 110 页起）再从规范场论一侧给出 't Hooft/赫克算符在分歧点处的定义。不妨把这三处看作三条线索各自的落点。
\end{yizhu}

基于我们的对偶提议，我们在第~\ref{sec:4}~节中对几何朗兰兹纲领在驯服分歧情形下应如何表述提出一个提议。（在数学上，基于 \cite{ref29} 以及上一段所引的其他结果，已有类似的提议。其阐述可参见综述 \cite{ref7} 第 9.4 节。某些特殊情形已在即将发表的工作 \cite{ref32} 中详加研究。）向野分歧的推广将在别处讨论 \cite{ref33}。

在第~\ref{sec:5}~节中，我们用规范场论来描述可作用于分歧点处的膜上的算符（它们是 \cite{ref8} 所研究的 't Hooft/赫克算符（Hecke operator）的推广）。这为第~\ref{sec:4}~节所处理的某些课题提供了一条更贴近具体的途径。在第~\ref{sec:6}~节中，我们借助靶空间为复余伴随轨道（complex coadjoint orbit）并配以超 Kähler 度量的西格玛模型，对分歧点处行为的某些方面给出更局部的描述。这类度量最初在 \cite{ref34,ref35} 中针对半单或幂零轨道（nilpotent orbit）构造，随后被推广到任意轨道 \cite{ref36,ref37}。这一点以及第~\ref{sec:3.8}~节的相关分析，是我们最接近于分析几何朗兰兹局部情形之处。最后，在第~\ref{sec:6}~节中我们还将简要描述本文所考虑类型面算符的一些弦理论构造。

\begin{yizhu}[物理图像补充]
余伴随轨道上的超 Kähler 度量有规范场论来源：它们可由纳姆方程（Nahm 方程）带边界条件的解得到。Kronheimer \cite{ref34,ref35} 借此构造出半单与幂零轨道的度量，Biquard 与 Kovalev \cite{ref36,ref37} 又把它推广到任意轨道。第~\ref{sec:3.8}~节（PDF 第 74 页起）用它解释 $\MH(\alpha,\beta,\gamma;p)$ 在 $(\alpha,\beta,\gamma)$ 非正则处的局部奇点：此时局部模型是复余伴随轨道，等价地（在另一个复结构下）是 $T^*(G/\mathbb{L})$，$\mathbb{L}$ 取遍列维子群。正是这套有限维模型，使分歧点附近的几何朗兰兹问题可以被具体分析。
\end{yizhu}

利用共形场论，人们已对允许任意阶极点的几何朗兰兹纲领的统一处理提出一个提议 \cite{ref6}。这项工作在 \cite{ref7} 中有综述。遗憾的是，本文既未与共形场论的使用建立联系，也未与这一统一表述建立联系。当然，我们希望最终能理解得更多。

\begin{yizhu}[含糊点拨]
\cite{ref6} 与 \cite{ref7} 的角色不同：前者是提议本身，题目为 \emph{Local Geometric Langlands Correspondence And Affine Kac--Moody Algebras}（math.RT/0508382），用的是顶点代数与仿射 Kac--Moody 代数的语言，因而可容纳任意阶极点；后者是综述，引言前文所引的正是其第 9.4 节。本文走四维面算符路线，目前只覆盖单极点即驯服分歧（第~\ref{sec:4}~节），故作者在此明言两条路线尚未打通。接口在局部：第~\ref{sec:6}~节用复余伴随轨道的西格玛模型逼近局部情形，并自承那是全文“最接近于”局部分析之处。
\end{yizhu}

附录 A 回顾了一些群论背景。附录 B 给出记号索引。本文用到的关于 Montonen--Olive 对偶、希钦模空间等大量事实，在 \cite{ref8} 中有更充分的解释。

我们感谢 A. Braverman、D. Gaitsgory、E. Frenkel 与 D. Kazhdan 耐心而极有助益的讲解。我们也感谢 J. Andersen、P. Aspinwall、M. F. Atiyah、D. Ben-Zvi、R. Bezrukavnikov、R. Bielawski、R. Dijkgraaf、R. Donagi、N. Hitchin、L. Jeffrey、A. Kapustin、P. Kronheimer、Y. Laszlo、H. Nakajima、C. Sorger 与 M. Thaddeus 等人提出的各种有益评论与建议。SG 的研究部分受 DOE 基金 DE-FG03-92-ER40701 资助。EW 的研究部分受 NSF 基金 PHY-0503584 资助。
